\chapter{Se repérer et représenter l'espace}
\label{corriges:18}
\section*{Exercice 18.1 -- Coordonnées et perspective}
$A$ est un sommet : il se trouve à l'extrémité $x=6$, sur la face $y=0$ et à la hauteur maximale $z=3$. $B$ appartient à la face $y=4$, avec $x=2$ et $z=1$. Les arêtes fuyantes sont représentées avec un angle et un coefficient choisis ; on ne les mesure pas comme si la perspective était une vue en vraie grandeur.

\section*{Exercice 18.2 -- Section d'une boule}
Le rayon de section vaut $\sqrt{13^2-5^2}=\sqrt{144}=12\unite{cm}$, par Pythagore. Son aire vaut $144\pi\unite{cm^2}$. À distance $13$, le plan est tangent et la section se réduit à un point. À distance $14>13$, il n'y a pas d'intersection.

\section*{Exercice 18.3 -- Section et patron d'un cône}
Le petit cône a un facteur de longueur $2/8=1/4$, donc un rayon de $6/4=1{,}5\unite{cm}$ et un facteur de volume $1/64$.
La génératrice du grand cône vaut $\sqrt{6^2+8^2}=10\unite{cm}$. Le secteur du patron a pour rayon $10$ et pour angle $360\times6/10=216\degre$, afin que son arc ait la longueur $2\pi6$.

\section*{Exercice 18.4 -- Sur la Terre sphérique}
Le point est dans l'hémisphère sud, à $20\degre$ au sud de l'équateur et sur le méridien situé à $40\degre$ à l'ouest de celui d'origine. L'arc de grand cercle mesure $40000\times18/360=2000\unite{km}$. Non : les parallèles hors équateur ont des rayons plus petits, donc leurs arcs de même angle sont plus courts. Un trajet le long d'un parallèle n'est pas généralement le plus court trajet sur la sphère.
